\documentclass[10pt]{article}
\usepackage[T1]{fontenc}
\usepackage[utf8]{inputenc}
\usepackage[french]{babel}
\usepackage{amsmath,amssymb,mathrsfs,mathpazo}
\usepackage{geometry}
\geometry{a4paper,margin=2cm}
\usepackage{graphicx,array,tabularx,tabularray,colortbl,multirow,multicol,booktabs,enumitem,pifont,ulem,fancybox,diagbox,cancel}
\usepackage{pstricks,pstricks-add,pst-plot,pst-tree,pst-node,pst-text,pst-func,pst-math,pst-bspline,pst-3dplot,pst-eucl,pst-solides3d}
\let\pscalculate\undefined
\usepackage{pst-bezier}
\usepackage{tikz,pgfplots,tkz-tab}
\usetikzlibrary{arrows,arrows.meta,patterns,calc,positioning,trees,decorations.markings,shapes.geometric,shapes.misc,shapes.arrows,intersections,decorations.pathmorphing,decorations.pathreplacing,backgrounds,shadows,babel}
\pgfplotsset{compat=1.18}
\usepgfplotslibrary{fillbetween}
\usepackage[np]{numprint}
\usepackage{listings,makecell,pas-tableur,twoopt}
\let\barre\undefined
\usepackage[pstricks]{tablvar}
\definecolor{darkgreen}{RGB}{0,100,0}
\providecommand{\red}{\color{red}}
\providecommand{\blue}{\color{blue}}
\providecommand{\green}{\color{green}}
\providecommand{\black}{\color{black}}
\providecommand{\white}{\color{white}}
\providecommand{\gray}{\color{gray}}
\providecommand{\cyan}{\color{cyan}}
\providecommand{\magenta}{\color{magenta}}
\providecommand{\yellow}{\color{yellow}}
\providecommand{\violet}{\color{violet}}
\providecommand{\orange}{\color{orange}}
\providecommand{\euro}{\texteuro}
\providecommand{\EUR}{\texteuro}
\providecommand{\ieme}{\textsuperscript{e}}
\providecommand{\degres}{\ensuremath{^\circ}}
\providecommand{\degree}{\ensuremath{^\circ}}
\providecommand{\vv}[1]{\overrightarrow{#1}}
\providecommand{\Coord}[1]{#1}
\providecommand{\decosix}{}
\providecommand{\starredbullet}{\ensuremath{\bullet}}
\pagestyle{empty}
\setlength{\parindent}{0pt}
\usepackage{adjustbox}\DeclareUnicodeCharacter{2236}{:}\DeclareUnicodeCharacter{2190}{\ensuremath{\leftarrow}}
\providecommand{\R}{\mathbb{R}}
\providecommand{\N}{\mathbb{N}}
\providecommand{\Z}{\mathbb{Z}}
\providecommand{\Q}{\mathbb{Q}}
\providecommand{\D}{\mathbb{D}}
\providecommand{\C}{\mathbb{C}}
\providecommand{\e}{{\mathrm{e}}}
\providecommand{\vect}[1]{\overrightarrow{#1}}
\providecommand{\vectt}[1]{\overrightarrow{\mathrm{#1}}}
\providecommand{\vvt}[1]{\overrightarrow{\mathrm{#1}}}
\providecommand{\np}[2][]{{#2~#1}}
\providecommand{\vv}[1]{{\overrightarrow{#1}}}
\renewcommand{\d}{{\mathrm{d}}}
\providecommand{\degres}{{^\circ}}
\providecommand{\degre}{{^\circ}}
\providecommand{\blue}{\color{blue}}
\providecommand{\red}{\color{red}}
\providecommand{\enskip}{\quad}
\providecommand{\llbracket}{\lbrack\!\lbrack}
\providecommand{\rrbracket}{\rbrack\!\rbrack}
\providecommand{\euro}{{\text{€}}}
\providecommand{\no}{{\text{n°}}}
\providecommand{\interval}[2]{{[#1;#2]}}
\providecommand{\intervFO}[2]{{[#1;#2[}}
\providecommand{\dotfill}{\ldots}

\providecommand{\ds}{\displaystyle}
\providecommand{\textding}[1]{#1}
\providecommand{\Oij}{\(\mathrm{O};\vec{\imath},\vec{\jmath}\)}
\providecommand{\Oijk}{\(\mathrm{O};\vec{\imath},\vec{\jmath},\vec{k}\)}
\providecommand{\Ouv}{\(\mathrm{O};\vec{u},\vec{v}\)}
\providecommand{\intervFF}[2]{\ensuremath{[#1;#2]}}
\renewcommand{\euro}{\ifmmode\text{\texteuro}\else\texteuro\fi}
\DeclareUnicodeCharacter{00B0}{\ifmmode^{\circ}\else\textdegree\fi}
\DeclareUnicodeCharacter{20AC}{\ifmmode\text{\texteuro}\else\texteuro\fi}
% Énoncés : Éducation nationale. Transcription et corrigés : APMEP, Denis Vergès et contributeurs. Découpage et classement : CoopMaths, Sébastien Lozano, Éric Elter et contributeurs. Adaptation de la présentation : Maths Explorer.
% CC BY-SA : https://coopmaths.fr/www/about/
% https://coopmaths.fr/alea/static/e3c/2020/tex/e3c_2020_00_sujet33_1_cor.tex
\begin{document}
\noindent\textit{APMEP / CoopMaths. CC BY-SA. Présentation adaptée par Maths Explorer.}\par\medskip
\subsection*{Question 1}

On sait qu'un vecteur directeur de la droite \((\textit{d})\) est \(\vec{u} \begin{pmatrix} -b \\ a \end{pmatrix}\), soit ici \(\vec{u} \begin{pmatrix} 3 \\ 2 \end{pmatrix}\).

\subsection*{Question 2}

Un vecteur directeur de la droite est \(\vec{u} \begin{pmatrix} 3 \\ 2 \end{pmatrix}\), donc un vecteur normal est par exemple \(\vec{v} \begin{pmatrix} -2 \\ 3 \end{pmatrix}\) ou encore \(-2 \vec{v} \begin{pmatrix} 1 \\ -\dfrac{3}{2} \end{pmatrix}\).

\subsection*{Question 3}

\(\overrightarrow{EB} = \overrightarrow{BA}\) entraîne en faisant intervenir \(\ C\) :
\(\overrightarrow{EC} + \overrightarrow{CB} = \overrightarrow{BC} + \overrightarrow{CA}\),

d'où : \(\overrightarrow{AC} = 2 \overrightarrow{BC} + \overrightarrow{CE}\).

De même : \(\overrightarrow{ED} = 2 \overrightarrow{BC}\) entraîne \(\overrightarrow{EC} + \overrightarrow{CD} = 2 \overrightarrow{BC}\),

d'où : \(\overrightarrow{CD} = 2 \overrightarrow{BC} + \overrightarrow{CE}\).

Conclusion : \(\overrightarrow{AC} = \overrightarrow{CD}\),

ce qui démontre que \(\ C \ \text{est le milieu de} \ [AD]\).

\subsection*{Question 4}

\(\vec{u}\) et \(\vec{v}\) sont orthogonaux, donc \(\vec{u} \cdot \vec{v} = 0\), soit :
\begin{align*}
&9(-x + 4) + 7(2x - 5) = 0 \\
\iff &-9x + 36 + 14x - 35 = 0 \\
\iff &5x + 1 = 0 \\
\iff &x = -\dfrac{1}{5}.
\end{align*}

\subsection*{Question 5}

Avec \(\overrightarrow{AC} \begin{pmatrix} 4 \\ 1 \end{pmatrix}\) et \(\overrightarrow{BD} \begin{pmatrix} -5 \\ 4 \end{pmatrix}\), on obtient :
\[
\overrightarrow{AC} \cdot \overrightarrow{BD} = 4 \times (-5) + 1 \times 4 = -20 + 4 = -16.
\]

\end{document}
